\documentclass[11pt,a4paper]{article}

% ======================================================
% LEBENSLAUF-VORLAGE (Sidebar-Format, Deutsch)
% ------------------------------------------------------
% So verwenden Sie diese Vorlage:
%   1. Alle Platzhalter in [eckigen Klammern] ersetzen.
%   2. Datumsangaben im Format MM/JJJJ eintragen
%      (laufende Taetigkeiten: MM/JJJJ--heute).
%   3. photo.jpg und signature.png in den Projektordner
%      legen. Fehlen die Dateien, kompiliert die Vorlage
%      trotzdem und zeigt einen Platzhalterrahmen.
%   4. Nicht benoetigte Eintraege/Abschnitte loeschen
%      oder mit % auskommentieren.
% ======================================================

% ======================================================
% PACKAGES
% ======================================================
\usepackage[a4paper,
left=7.2cm,
right=1.5cm,
top=1.8cm,
bottom=2cm]{geometry}

\usepackage[T1]{fontenc}
\usepackage[utf8]{inputenc}
\usepackage[english,ngerman]{babel}

\usepackage{changepage}
\usepackage{graphicx}
\usepackage{xcolor}
\usepackage{tikz}
\usepackage{eso-pic}
\usepackage{tabularx}
\usepackage{array}
\usepackage{titlesec}
\usepackage[hidelinks]{hyperref}
\usepackage{fontawesome5}
\usepackage{parskip}
\usepackage{enumitem}
\usepackage{ifthen}
\hypersetup{
    colorlinks=true,
    linkcolor=blue,
    urlcolor=blue
}

% ======================================================
% COLORS
% ======================================================
\definecolor{sidebarblue}{RGB}{220,232,245}

% ======================================================
% PHOTO AND SIGNATURE FILES
% (Dateinamen bei Bedarf hier anpassen)
% ======================================================
\newcommand{\cvphoto}{photo.jpg}
\newcommand{\cvsignature}{signature.png}

% ======================================================
% SIDEBAR ONLY ON FIRST PAGE
% ======================================================
\AddToShipoutPictureBG*{
\begin{tikzpicture}[remember picture,overlay]

\fill[sidebarblue]
(current page.north west)
rectangle ([xshift=6.5cm]current page.south west);

\end{tikzpicture}
}

% ======================================================
% SECTION FORMAT
% ======================================================
\titleformat{\section}
{\Large\bfseries}
{}
{0pt}
{}
[\titlerule]

\titlespacing*{\section}
{0pt}{12pt}{8pt}

% ======================================================
% ENTRY COMMANDS
% ======================================================
% \cventry{Zeitraum}{Titel/Position}{Institution/Firma}{Beschreibung}
\newcommand{\cventry}[4]
{
\vspace{0.8em}

\begin{tabularx}{\textwidth}{p{3.5cm}X}
\textbf{#1}
&
\textbf{#2}
\\[0.2em]
&
\textit{#3}
\\[0.5em]
&
\begin{minipage}[t]{0.95\linewidth}
#4
\end{minipage}
\end{tabularx}

\vspace{0.1em}
}

% \simpleentry{Jahr}{Beschreibung} — z.B. fuer Auszeichnungen
\newcommand{\simpleentry}[2]
{
\begin{tabularx}{\textwidth}{p{3cm}X}
\textbf{#1} & #2
\end{tabularx}

\vspace{0.2em}
}

% \skillentry{Kategorie}{kommagetrennte Liste}
\newcommand{\skillentry}[2]{%
    \begin{tabularx}{\textwidth}{@{} p{5cm} X @{}}
        \textbf{#1} & #2 
    \end{tabularx}\par
}

% Beendet das Sidebar-Layout und schaltet auf normale
% Seitenraender um (ab Seite 2 aufrufen).
\newcommand{\switchtonormalpages}{
\vfill
\clearpage
\newgeometry{
left=2cm,
right=1.5cm,
top=1.8cm,
bottom=2cm
}
}

\setlist[itemize]{
leftmargin=1.2em,
topsep=2pt,
itemsep=1pt,
parsep=0pt,
partopsep=0pt
}

% ======================================================
% DOCUMENT
% ======================================================
\begin{document}

% ======================================================
% SIDEBAR CONTENT
% ======================================================
\begin{tikzpicture}[remember picture,overlay]

\node[
anchor=north west,
inner sep=0pt
]
at ([xshift=0.6cm,yshift=-0.8cm]current page.north west)
{

\begin{minipage}[t]{5.2cm}

\centering
% ======================================================
% PHOTO
% ======================================================

\centering
\IfFileExists{\cvphoto}{%
    \includegraphics[width=4cm,height=5cm]{\cvphoto}%
}{%
    \fbox{\parbox[c][5cm][c]{3.8cm}{\centering Foto\\einf\"ugen\\\texttt{(photo.jpg)}}}%
}

\vspace{0.6cm}

% ======================================================
% NAME
% ======================================================

{\fontsize{22}{24}\selectfont\bfseries [Vorname]}\\
\vspace{0.2cm}
{\fontsize{22}{24}\selectfont\bfseries [Nachname]}

\vspace{0.3cm}

{\fontsize{15}{18}\selectfont\bfseries Lebenslauf}

\raggedright
\vspace{0.5cm}
% ======================================================
% PERSONAL DATA
% ======================================================

{\large\bfseries Pers\"onliche Daten}

\vspace{0.4cm}

\textbf{Geburtsdatum}\\
TT.MM.JJJJ

\vspace{0.3cm}

% Optionale Felder — bei Bedarf einkommentieren:

% \textbf{Geburtsort}\\
% [Ort, Land]

% \vspace{0.3cm}

% \textbf{Staatsangeh\"origkeit}\\
% [Staatsangeh\"origkeit]

% \vspace{0.3cm}

% \textbf{Familienstand}\\
% [ledig/verheiratet]

\vspace{0.8cm}

% ======================================================
% CONTACT
% ======================================================

{\large\bfseries Kontakt}

\vspace{0.4cm}

\faPhone\ +49 XXX XXXXXXXX

\vspace{0.25cm}

\faEnvelope\ \href{mailto:vorname.nachname@example.com}{vorname.nachname@example.com}

\vspace{0.25cm}

\faMapMarker*\ [Stadt, Land]

\vspace{0.25cm}

\faLinkedin\ \href{https://linkedin.com/in/benutzername}{linkedin.com/in/benutzername}

\vspace{0.25cm}

% \faGithub\ \href{https://github.com/benutzername}{github.com/benutzername}

\vspace{0.8cm}

% ======================================================
% LANGUAGES
% ======================================================

{\large\bfseries Sprachkenntnisse}

\vspace{0.4cm}

[Sprache 1] — [Niveau, z.B. Muttersprache]

\vspace{0.2cm}

[Sprache 2] — [Niveau, z.B. C1]

\vspace{0.2cm}

[Sprache 3] — [Niveau, z.B. B1]

\end{minipage}

};

\end{tikzpicture}

% ======================================================
% MAIN CONTENT
% ======================================================

% \section{Kurzprofil}

% [2--3 Saetze: Wer sind Sie, was koennen Sie, was suchen Sie?]

% ======================================================
% 1. Education
% ======================================================

\section{Ausbildung}

\cventry
{MM/JJJJ--heute}
{[Abschluss und Studiengang, z.B. M.Sc. ...]}
{[Hochschule, Ort]}
{Note: [X{,}X] % oder ECTS-Schnitt; Zeile loeschen, falls nicht gewuenscht
\\[0.2em]
Schwerpunkte: [Schwerpunkt 1], [Schwerpunkt 2], [Schwerpunkt 3]}

\cventry
{MM/JJJJ--MM/JJJJ}
{[Abschluss und Studiengang, z.B. B.Sc. ...]}
{[Hochschule, Ort]}
{
Note: [X{,}X]

\vspace{0.4em}
\textbf{Abschlussarbeit} (MM/JJJJ--MM/JJJJ): \textit{[Titel der Arbeit]}
\begin{itemize}
\item {[Kernaufgabe oder wichtigstes Ergebnis, moeglichst mit Zahl]}
\end{itemize}
}

% ======================================================

\section{Berufserfahrung}

\cventry
{MM/JJJJ--heute}
{[Position, z.B. Werkstudent/in]}
{[Firma, Ort]}
{
\begin{itemize}
\item {[Taetigkeit mit messbarem Ergebnis: Was? Womit? Resultat?]}
\item {[Weitere Taetigkeit oder Verantwortung]}
\end{itemize}
}

\cventry
{MM/JJJJ--MM/JJJJ}
{[Position]}
{[Firma, Ort]}
{
\begin{itemize}
\item {[Taetigkeit mit messbarem Ergebnis]}
\end{itemize}
}

% ======================================================
\switchtonormalpages

\section{Projekte}

\cventry
{MM/JJJJ--MM/JJJJ}
{[Projekttitel]}
{[Institution oder Kontext]}
{
\begin{itemize}
\item {[Was wurde gemacht?]}
\item {[Welche Methoden/Werkzeuge?]}
\item {[Ergebnis, moeglichst mit Zahl]}
\end{itemize}
}

\cventry
{MM/JJJJ--MM/JJJJ}
{[Projekttitel]}
{[Institution oder Kontext]}
{
\begin{itemize}
\item {[Kurzbeschreibung mit Ergebnis]}
\end{itemize}
}

% ======================================================

% \section{Auszeichnungen \& Erfolge}

% \simpleentry
% {JJJJ}
% {[Auszeichnung, Wettbewerb oder Stipendium mit Einordnung]}

% ======================================================

\section{Kenntnisse \& F\"ahigkeiten}

\skillentry
{Programmiersprachen}
{[Sprache 1], [Sprache 2], [Sprache 3]}

\skillentry
{Simulation}
{[Methode/Tool 1], [Methode/Tool 2]}

\skillentry
{Software \& Werkzeuge}
{[Tool 1], [Tool 2], [Tool 3]}

\skillentry
{Labor- und Forschungskenntnisse}
{[Kenntnis 1], [Kenntnis 2], [Kenntnis 3]}

\skillentry
{Laborger\"ate}
{[Ger\"at 1], [Ger\"at 2], [Ger\"at 3]}

% ======================================================

\section{Hobbys \& Interessen}

[Hobby 1], [Hobby 2], [Hobby 3]

\vfill

% ======================================================
% PLACE, DATE AND SIGNATURE
% ======================================================
\vspace{0.4cm}
[Stadt], \today

\IfFileExists{\cvsignature}{%
    \includegraphics[height=2cm]{\cvsignature}%
}{%
    \vspace{2cm}%
}

\vspace{-0.5cm}

\rule{5cm}{0.4pt}

[Vorname Nachname]

\end{document}